\section{Preliminaries}
\label{sec:preliminaries}

This section fixes notation and the formal interfaces used by the construction: target-shaped commitments,
the BB-tran rational hash, the imported DKLW signature layer, non-interactive zero-knowledge arguments,
issuance and ARC security notions, and the PRF.

\subsection{Notation}
Let \(\lambda\) denote the security parameter and let \(\negl(\lambda)\) be any positive function negligible in
\(\lambda\).  We write PPT for probabilistic polynomial time.  Vectors are written in bold lower-case letters and
matrices in bold upper-case letters.

We work over the cyclotomic ring
\[
  \R=\Z[X]/(X^N+1)
\]
for power-of-two \(N\), and its quotient \(\Rq=\R/q\R\).  We identify \(\Rq\) with polynomials of degree
\(<N\) with coefficients in \(\Z_q\).  Whenever a field element is interpreted as a signed integer, we use its
signed representative in \([-q/2,q/2]\).  For \(x\in\Rq\), the notation \(\|x\|_\infty\le B\) means that every
coefficient of \(x\) lies in \([-B,B]\); the same convention is used component-wise for vectors over \(\Rq\).

The semantic signed message is
\[
  M:=\vecm\Vert\veck\in\Rq^{\ell_M},
\]
where \(\vecm\) encodes holder attributes and \(\veck\) is the secret PRF key used for rate limiting.  The first
target parameter set uses \(\ell_M=1\), so the encoding of \(\vecm\Vert\veck\) into one ring polynomial is an
explicit parameter assumption.  If a deployment needs a longer semantic message, \(\ell_M\) must be increased and
all commitment and SmallWood row counts recomputed.

The PRF key is sampled from a finite key space
\(\mathcal K_{\mathsf{key}}\subseteq \Fq^{\ell_{\mathsf{key}}}\) fixed by the public parameters.  In the bounded-key
instantiation, we require
\[
  \mathcal K_{\mathsf{key}}\subseteq
  \{\veck\in\Fq^{\ell_{\mathsf{key}}}:\|\veck\|_\infty\le \BoundMsg\},
\]
and the PRF assumption is made for a uniform key from this same key space.  The public coefficient bound for
holder attributes and commitment randomness is denoted \(\BoundMsg\), and the short-preimage bound for signatures
is denoted \(\BoundSig\).

Whenever the CRT decomposition \(\Rq\cong\prod_i F_i\) is used, we write
\[
  (x_i)_i\Had(y_i)_i := (x_i y_i)_i,
  \qquad
  (x_i)_i\Hdiv (y_i)_i := (x_i y_i^{-1})_i
\]
whenever every \(y_i\neq 0\).  Thus \(\Had\) and \(\Hdiv\) denote slotwise multiplication and slotwise division in
CRT coordinates.

\subsection{Commitments and target-shaped MLWE hiding}
\label{subsec:prelim-commit}
\label{subsec:commitments}

\begin{definition}[Commitment scheme]\label{def:commitment-scheme}
A non-interactive commitment scheme is a tuple of PPT algorithms
\[
  \ComScheme=(\ComSetup,\Commit,\ComVerify)
\]
with the following syntax.
\begin{itemize}[leftmargin=1.25em]
  \item \(\ComPublicParam\gets \ComSetup(\SecParam,\PublicParamsGen)\) outputs public commitment parameters.
  \item \((c,d)\gets \Commit(m)\) takes a message \(m\) and outputs a commitment \(c\) and an opening \(d\).
  \item \(\ComVerify(c,m,d)\in\{0,1\}\) deterministically checks whether \(d\) is a valid opening of \(c\) to \(m\).
\end{itemize}
The scheme is correct if honestly generated commitments always verify.  It is computationally binding if no PPT
adversary can open one commitment to two distinct messages except with negligible probability.  It is hiding if
commitments to any two messages are computationally indistinguishable.
\end{definition}

The issuance protocol uses a compact target-shaped lattice commitment to the semantic message.  The shape is a
local specialization of the Ajtai/BDLOP/ABDLOP commitment lineage used in lattice zero-knowledge
protocols~\cite{BDLOP18,LNP22}.  The source precedent is a lattice commitment with hiding based on
Module-LWE-type pseudorandomness and binding based on Module-SIS-type short-relation hardness.  The exact compact
one-target form below is the ARC-SPRUCE specialization.

Let
\[
  M\in\Rq^{\ell_M},\qquad s\in\R^{k_s},\qquad e\in\R^{n_c}.
\]
The public matrices are
\[
  C_M\in\Rq^{n_c\times \ell_M},
  \qquad
  A_s\in\Rq^{n_c\times k_s},
\]
and the commitment is
\begin{equation}
  c=C_M M+A_s s+e\in\Rq^{n_c}.
  \label{eq:target-commitment}
\end{equation}
Equivalently, with
\[
  C=[C_M\mid A_s\mid I_{n_c}],
\]
this is the block equation \(c=C[M;s;e]\).  The identity block is intentional: no random invertible error block is
required for the compact commitment used here.

\begin{definition}[Opening relation for the target-shaped commitment]\label{def:target-open}
\(\Open(c;M,s,e)=1\) if and only if \eqref{eq:target-commitment} holds in \(\Rq^{n_c}\), \(M\) is a valid
encoding of the semantic message, the key-space constraints inside \(M\) hold, and the coefficients of opened values
satisfy the prescribed bounds.  For the target set used in Section~\ref{sec:parameters}, \(\ell_M=1\), \(k_s=2\),
\(n_c=1\), and the coefficients of \(M,s,e\) are bounded by \(\BoundMsg=8\) unless stated otherwise.
\end{definition}

\begin{definition}[Module-LWE hiding assumption]\label{ass:mlwe-hiding}
For parameters \((\R,q,n_c,k_s,\chi_s,\chi_e)\), define
\[
  \Adv^{\mathsf{mlwe\mbox{-}hide}}_{\Adv}(\lambda)
  :=
  \left|
    \Pr\bigl[\Adv(A_s,A_s s+e)=1\bigr]
    -
    \Pr\bigl[\Adv(A_s,U)=1\bigr]
  \right|,
\]
where \(A_s\getsunif\Rq^{n_c\times k_s}\), \(s\leftarrow\chi_s^{k_s}\), \(e\leftarrow\chi_e^{n_c}\), and
\(U\getsunif\Rq^{n_c}\).  The assumption
\(\MLWE\text{-}\mathsf{Hiding}_{\R,q,n_c,k_s,\chi_s,\chi_e}\) states that
\(\Adv^{\mathsf{mlwe\mbox{-}hide}}_{\Adv}(\lambda)\le\negl(\lambda)\) for every PPT adversary \(\Adv\).
\end{definition}

\begin{lemma}[Computational hiding of the target-shaped commitment]\label{lem:commit-hiding}
Under Assumption~\ref{ass:mlwe-hiding}, commitments to any two fixed semantic messages \(M_0,M_1\) are
computationally indistinguishable.
\end{lemma}

\begin{proof}
For fixed \(M\), the distribution of \(c=C_MM+A_s s+e\) is a deterministic offset \(C_MM\) plus
\(A_s s+e\).  By Assumption~\ref{ass:mlwe-hiding}, this is computationally indistinguishable from the same offset
plus a uniform \(U\).  Since \(U+C_MM\) is uniform over \(\Rq^{n_c}\), the resulting distribution is independent of
\(M\).  Applying the argument to \(M_0\) and \(M_1\) gives the claim.
\end{proof}

\begin{definition}[Module-SIS binding assumption]\label{ass:msis-binding}
Let
\[
  L_{\mathsf{bind}}=\ell_M+k_s+n_c.
\]
For uniform \(C_M\getsunif\Rq^{n_c\times\ell_M}\) and \(A_s\getsunif\Rq^{n_c\times k_s}\), define
\[
  \Adv^{\mathsf{msis\mbox{-}bind}}_{\Adv}(\lambda)
  :=
  \Pr\!\bigl[
    y\neq 0\ \wedge\ [C_M\mid A_s\mid I_{n_c}]y=0\bmod q\ \wedge\
    \|y\|_\infty\le \beta_{\mathsf{bind}}
  \bigr],
\]
where \(y\in\R^{L_{\mathsf{bind}}}\) is output by \(\Adv(C_M,A_s)\).  The assumption
\(\MSIS\text{-}\mathsf{Bind}_{\R,q,n_c,L_{\mathsf{bind}},\beta_{\mathsf{bind}}}\) states that
\(\Adv^{\mathsf{msis\mbox{-}bind}}_{\Adv}(\lambda)\le\negl(\lambda)\) for every PPT adversary \(\Adv\).
\end{definition}

\begin{lemma}[Binding of the target-shaped commitment]\label{lem:commit-binding}
If the same commitment \(c\) has two distinct bounded openings \((M,s,e)\) and \((M',s',e')\), then their
explicit difference gives a short non-zero solution to the Module-SIS instance of Assumption~\ref{ass:msis-binding}.
If all opened coefficients are bounded by \(\BoundMsg\), then the difference has coefficient bound
\(2\BoundMsg\) for those opened coordinates.
\end{lemma}

\begin{proof}
Subtracting the two equations \(c=C_MM+A_s s+e=C_MM'+A_s s'+e'\) yields
\[
  C_M(M-M')+A_s(s-s')+(e-e')=0.
\]
Since the openings are distinct, the block vector \((M-M',s-s',e-e')\) is non-zero.  Its coefficient bound is at
most twice the opening bound, and it is a Module-SIS solution for \([C_M\mid A_s\mid I_{n_c}]\).
\end{proof}

\subsection{Lattice trapdoors and trapdoor rational hashes}
\label{subsec:prelim-vsis}

\begin{definition}[Admissible lattice trapdoor suite]\label{def:trapdoor-suite}
A lattice trapdoor suite over \(\Rq\) consists of PPT algorithms
\[
  (\Setup,\TrapGen,\SampPre)
\]
together with a public bound \(\BoundSig\).  We say that the suite is admissible for the parameters used in this
paper if the following interface holds.
\begin{itemize}[leftmargin=1.25em]
  \item \(\PublicParamsTrapdoor\gets\Setup(\SecParam,\PublicParamsGen)\) outputs dimensions \((n,m)\), modulus
  \(q\), sampler parameters, and the accepted shortness bound \(\BoundSig\).
  \item \((\mA,\trapdoor)\gets\TrapGen(\PublicParamsTrapdoor)\) outputs a public matrix
  \(\mA\in\Rq^{n\times m}\) together with a trapdoor \(\trapdoor\), and the distribution of \(\mA\) is
  statistically or computationally close to uniform over \(\Rq^{n\times m}\) as required by the invoked
  signature theorem.
  \item On input \((\PublicParamsTrapdoor,\trapdoor,T)\) with \(T\in\Rq^n\), the sampler
  \(\SampPre(\trapdoor,T)\) outputs \(\vecu\in\R^m\) such that \(\mA\vecu=T\) in \(\Rq^n\), except with
  negligible failure probability, and \(\|\vecu\|_\infty\le \BoundSig\) except with negligible probability.
  \item The distributional property required by the DKLW hash-and-preimage signature theorem holds for the
  produced preimages at the chosen Gaussian width and public bound.
\end{itemize}
This definition is an interface assumption for the signature-layer import.  Appendix~\ref{app:gaussian-sampler}
describes the NTRU/Antrag-based candidate used for parameter tuning, but Antrag alone is not cited here as a proof
of this full admissible-suite interface.
\end{definition}

\begin{definition}[Trapdoor rational hash]\label{def:trapdoor-rational-hash}
A trapdoor rational hash scheme over \(\Rq\) is a tuple
\[
  \TDRH=(\Setup,\TrapGen,\Eval,\SampPre)
\]
consisting of the lattice trapdoor suite together with a public evaluation algorithm.  The setup output
\(\PublicParamsTrapdoor\) contains a public hash description \(B\), a characteristic space
\(\mathcal M_{\mathsf{sig}}\), a randomness space \(\mathcal X_{\mathsf{sig}}\), and an admissibility error bound
\(\delta\).  For fixed \(B\), the algorithm
\[
  \Eval(\PublicParamsTrapdoor,\muSig,\chiSig)
\]
takes \(\muSig\in\mathcal M_{\mathsf{sig}}\) and \(\chiSig\in\mathcal X_{\mathsf{sig}}\) and outputs either a
target \(T\in\Rq^n\) or \(\bot\) if the input is inadmissible.  Define the admissible characteristic subspace
\[
  \mathcal M^{\mathsf{sig}}_B:=
  \left\{
    \muSig\in\mathcal M_{\mathsf{sig}}:
    \Pr_{\chiSig\gets\mathcal X_{\mathsf{sig}}}
    \bigl[\Eval(\PublicParamsTrapdoor,\muSig,\chiSig)=\bot\bigr]
    \le \delta
  \right\}.
\]
Correctness requires that for every \(\muSig\in\mathcal M^{\mathsf{sig}}_B\) and every admissible
\(\chiSig\) such that \(\Eval(\PublicParamsTrapdoor,\muSig,\chiSig)=T\neq\bot\), the output
\(\vecu\gets\SampPre(\trapdoor,T)\) satisfies
\[
  \mA\vecu=T
  \qquad\text{and}\qquad
  \|\vecu\|_\infty\le \BoundSig
\]
except with negligible probability.
\end{definition}

\paragraph{BB-tran translation.}
DKLW's BB-tran row has message or characteristic \(u\), randomness \((x_0,x_1)\), and rational function
\[
  (u^\top,x_0^\top)\tilde b_0+\frac{1}{\tilde b_1-x_1}.
\]
In the committed-message construction, the rational-hash characteristic is issuer sampled and is not the semantic
message.  We use
\[
  u\leftrightarrow \muSig,
  \qquad
  x_0\leftrightarrow x_0,
  \qquad
  x_1\leftrightarrow x_1.
\]
The public description is \(B=(B_0,B_1,B_2,B_3)\), with \(B_0,B_3\in\Rq^n\) and public linear maps
\(B_1,B_2\) into \(\Rq^n\).  With
\[
  \chiSig=(x_0,x_1),
\]
define
\begin{equation}
  h_{\muSig,\chiSig}(B):=B_0+B_1\muSig+B_2x_0+Z,
  \qquad
  Z=(B_3-\mathbf 1_nx_1)^{-1}.
  \label{eq:htran-target}
\end{equation}
The input is admissible if \(B_3-\mathbf 1_nx_1\in\Rq^{\times n}\), equivalently if every CRT slot of every
coordinate is non-zero.  The witness form is
\begin{equation}
  (B_3-\mathbf 1_nx_1)\Had Z=1_n,
  \qquad
  T=B_0+B_1\muSig+B_2x_0+Z+C_MM+A_s s+e.
  \label{eq:htran-witness-form}
\end{equation}
The semantic message \(M\) appears only through the commitment term.

\begin{lemma}[Eliminated form under admissibility]\label{lem:hash-cleared-admissible}
Let \(B_3-\mathbf 1_nx_1=(d_i)_i\) and
\(T-B_0-B_1\muSig-B_2x_0-C_MM-A_s s-e=(z_i)_i\) in CRT coordinates.  If \(d_i\neq 0\) for all CRT slots, then
\[
  T=h_{\muSig,\chiSig}(B)+C_MM+A_s s+e
\]
if and only if
\[
  (B_3-\mathbf 1_nx_1)\Had
  \bigl(T-B_0-B_1\muSig-B_2x_0-C_MM-A_s s-e\bigr)=1_n.
\]
\end{lemma}

\begin{proof}
In each CRT slot, the statement is \(z_i=d_i^{-1}\), which is equivalent to \(d_i z_i=1\) when \(d_i\neq0\).
Reassembling the slots proves the equivalence.
\end{proof}

\begin{definition}[(Interactive) \(\GenISIS_f\) interface]\label{def:genisis-interface}
Let
\[
  f:K\times\mathcal M_{\mathsf{sig}}\times\mathcal X_{\mathsf{sig}}\to\Rq^n
\]
be a keyed probabilistic family.  In the non-interactive \(\GenISIS_f\) interface, an adversary receives public
parameters and must output a fresh short preimage of a target of the form
\(f(\kappa,\muSig,\chiSig)\).  In the interactive \(\IntGenISIS_f\) interface, the adversary may obtain short
preimages of targets augmented by a linear committed-message term.  In the ARC-SPRUCE committed-message issuance
layer, this term is specialized locally to
\[
  C_MM+A_s s+e.
\]
\end{definition}

\subsection{Signatures and imported non-blind security}
\label{subsec:prelim-signatures}
The issuance layer builds on the non-blind vSIS hash-and-preimage signature layer from DKLW.  We record the
standard syntax and the conditional import used by this paper.

\begin{definition}[Signature scheme]\label{def:signature-scheme}
A signature scheme for a message space \(\mathcal M\) is a tuple of PPT algorithms
\[
  \Sigma=(\Setup,\KeyGen,\Sign,\SigVerify).
\]
It is correct if there exists a negligible function such that, for every \(\mu\in\mathcal M\),
\[
\Pr\!\left[
\begin{array}{l}
  \PP\gets\Setup(\SecParam);\\
  (\pk,\sk)\gets\KeyGen(\PP);\\
  \sig\gets\Sign(\sk,\mu);\\
  \SigVerify(\pk,\mu,\sig)=1
\end{array}
\right]
\ge 1-\negl(\lambda).
\]
\end{definition}

\begin{definition}[Strong existential unforgeability]\label{def:seuf-defs}
Let \(\Sigma\) be a signature scheme.  We use the standard notions of strong existential unforgeability under
chosen-message attack (sEUF-CMA), selective-message attack (sEUF-SMA), and strong selective unforgeability under
selective-message attack (sSUF-SMA), as in the DKLW formalization.  In each case a successful adversary outputs a
valid pair \((\mu^\star,\sig^\star)\) that was not returned by the relevant signing oracle, with the stronger notions
also forbidding a new valid signature on a previously signed message.
\end{definition}

\begin{definition}[vSIS-based hash-and-preimage signature]\label{def:vsis-template}
Let \(\TDRH\) be a trapdoor rational hash scheme with characteristic space \(\mathcal M_{\mathsf{sig}}\), randomness
space \(\mathcal X_{\mathsf{sig}}\), admissibility error \(\delta\), and public bound \(\BoundSig\).  The associated
vSIS signature scheme signs a DKLW characteristic \(\muSig\) by sampling \(\chiSig\leftarrow\mathcal X_{\mathsf{sig}}\),
computing \(T=\Eval(\PublicParamsTrapdoor,\muSig,\chiSig)\), and sampling
\(\vecu\gets\SampPre(\trapdoor,T)\).  Verification checks \(\chiSig\in\mathcal X_{\mathsf{sig}}\), admissibility,
\(\|\vecu\|_\infty\le\BoundSig\), and \(\mA\vecu=T\).
\end{definition}

\begin{theorem}[Specialized DKLW security import]\label{thm:dkwl-seuf-sma}
Let \(\Sigma_{\mathrm{tran}}\) be the non-blind vSIS signature scheme obtained from Definition~\ref{def:vsis-template}
by instantiating the rational family with the BB-tran form.  Assume the admissible trapdoor-suite interface of
Definition~\ref{def:trapdoor-suite}, the DKLW predicate and denominator-bound conditions, the required
strong-linear-independence condition for the instantiated rational functions, the required norm-compatibility
conversion between DKLW's norm convention and the public bound \(\BoundSig\), negligible inadmissibility errors,
and the corresponding strong hinted-vSIS hypothesis.  Then \(\Sigma_{\mathrm{tran}}\) is strongly existentially
unforgeable under selective-message attack.
\end{theorem}

\begin{proof}
This is the BB-tran specialization of the DKLW selective-query security theorem for their generic vSIS-based
signature family.  The statement here is conditional on the DKLW side conditions being satisfied by the concrete
ARC-SPRUCE parameters.
\end{proof}

\begin{corollary}[Conditional imported route to sEUF-CMA]\label{cor:dkwl-seuf-cma}
Suppose, in addition, that the exact hypotheses of the DKLW \(\GenISIS_f/\IntGenISIS_f\) lifting theorem for
BB-tran are met, including the required power-of-two cyclotomic setting, the modulus and message-bound condition
\(q/2>\gamma\), the dimension condition \(m=n\log q+\omega(\log\lambda)\), the smoothing lower bound on the
Gaussian parameter, and the theorem's stated bound loss.  Then the same non-blind signature layer obtains the
corresponding strong existential unforgeability guarantee under adaptive chosen-message attack.
\end{corollary}

\begin{proof}
This is a conditional import of DKLW's GenISIS/IntGenISIS lifting theorem.  It concerns only the non-blind
hash-and-preimage signature layer and does not imply one-more unforgeability of the interactive issuance protocol.
\end{proof}

\begin{remark}[Scope of the import]\label{rem:dkwl-scope}
Theorems~\ref{thm:dkwl-seuf-sma} and~\ref{cor:dkwl-seuf-cma} concern the non-blind hash-and-preimage signature
layer only.  DKLW's \(\IntGenISIS_f\) experiment includes a linear committed-message term; ARC-SPRUCE uses that
shape for issuance, but this paper does not claim a completed proof of one-more unforgeability for the interactive
issuance protocol.
\end{remark}

\subsection{Non-interactive zero-knowledge arguments of knowledge}
\label{subsec:prelim-nizk}
\label{subsec:nizk-ark}
We use the NIZK-ARK syntax of the SmallWood framework instantiated in the random-oracle model.

\begin{definition}[NIZKs in the random-oracle model]\label{def:nizk-ark}
Let \(\mathcal R\subseteq \mathcal X\times \mathcal W\) be a family of relations.  A transparent non-interactive
zero-knowledge argument of knowledge for \(\mathcal R\) in the random-oracle model, with oracle \(\RO\), is a
pair of algorithms \((\ZKProve^{\RO},\ZKVerify^{\RO})\):
\begin{itemize}[leftmargin=1.25em]
  \item \(\ZKProve^{\RO}(w\mid \mathcal R,x)\to \pi\) is probabilistic and, on input a statement-witness pair
  \((x,w)\in\mathcal R\), outputs a proof \(\pi\).
  \item \(\ZKVerify^{\RO}(\pi\mid \mathcal R,x)\to\{0,1\}\) is deterministic.
\end{itemize}
It satisfies the following properties.

\paragraph{Completeness.}
For every \((x,w)\in\mathcal R\),
\[
  \Pr\bigl[\ZKVerify^{\RO}(\pi\mid \mathcal R,x)=1 :
  \pi\gets\ZKProve^{\RO}(w\mid \mathcal R,x)\bigr]=1.
\]

\paragraph{Straightline-extractable knowledge soundness.}
There exists a PPT extractor \(\Ext\) such that for every PPT adversary \(\Adv\),
\[
  \Pr\Bigl[
    \ZKVerify^{\RO}(\pi\mid \mathcal R,x)=1\ \wedge\ (x,w)\notin \mathcal R
  \Bigr]\le \varepsilon,
\]
where \((\pi,Q)\leftarrow \Adv^{\RO}(x)\), \(Q\) is the set of random-oracle query/answer pairs made by
\(\Adv\), and \(w\leftarrow\Ext(Q,\pi)\).

\paragraph{Zero knowledge.}
There exists a PPT simulator \(\Sim\), which may program the random oracle, such that for every
\((x,w)\in\mathcal R\) and every PPT distinguisher \(\Adv\),
\[
  \Bigl|\Pr[\Adv^{\RO}(\pi_{\mathrm{real}})=1]
  -\Pr[\Adv^{\RO}(\pi_{\mathrm{sim}})=1]\Bigr|\le \xi,
\]
where
\[
  \pi_{\mathrm{real}}\gets \ZKProve^{\RO}(w\mid \mathcal R,x),
  \qquad
  \pi_{\mathrm{sim}}\gets \Sim^{\RO}(x).
\]
\end{definition}

We instantiate these proofs through the SmallWood PACS/PIOP framework compiled with Fiat--Shamir in the
random-oracle model.

\subsection{Issuance syntax and security}
\label{subsec:blind-sig-defs}
We keep the blind-signature-style syntax, but the protocol is stated as a committed-message issuance layer.

\begin{definition}[Issuance syntax]\label{def:bs-syntax}
The issuance layer is a tuple
\[
  \BlindSignature=(\Setup,\IKeyGen,\Issue,\Verify)
\]
where \(\Setup\) generates public parameters, \(\IKeyGen\) generates issuer keys \((\vk,\sk)\), \(\Issue\) is an
interactive protocol between issuer and holder, and \(\Verify(\vk,M,\sig)\) checks raw extracted witnesses.
\end{definition}

\begin{definition}[Issuance correctness]\label{def:bs-correctness}
The issuance layer is correct if there exists a negligible function such that, for every semantic message \(M\),
\[
\Pr\!\left[
\begin{array}{l}
  \PublicParamsBS\gets \Setup(\SecParam);\\
  (\vk,\sk)\gets \IKeyGen(\PublicParamsBS);\\
  \sig\gets \Holder^{\Issue:\langle\Issuer(\sk):\cdot\rangle}(M,\vk);\\
  \Verify(\vk,M,\sig)=1
\end{array}
\right]
\ge 1-\negl(\lambda).
\]
The probability is over setup, key generation, the honest holder and issuer coins, and the trapdoor sampler.
\end{definition}

\begin{definition}[Issuer-view hiding]\label{def:bs-blindness}
The issuance layer is issuer-view hiding if, for every PPT distinguisher \(\Adv\), the following left-right advantage
is negligible.  The challenger samples \(\PublicParamsBS\gets\Setup(\SecParam)\) and
\((\vk,\sk)\gets\IKeyGen(\PublicParamsBS)\).  The adversary outputs two equal-length semantic messages
\(M_0,M_1\).  The challenger samples \(b\getsunif\{0,1\}\), runs one honest holder issuance on input \(M_b\)
against the honest issuer, records the issuer's complete view \(\mathsf{View}_{\Issuer,b}\), and gives this view to
\(\Adv\).  If \(\Adv\) outputs \(\hat b\), then
\[
  \Adv^{\mathsf{ivh}}_{\BlindSignature,\Adv}(\lambda)
  :=
  \left|\Pr[\hat b=b]-\frac12\right|.
\]
The protocol is issuer-view hiding if \(\Adv^{\mathsf{ivh}}_{\BlindSignature,\Adv}(\lambda)\le\negl(\lambda)\) for every
PPT \(\Adv\).
\end{definition}

\begin{definition}[One-more unforgeability]\label{def:bs-one-more}
For \(Q\in\mathbb N\), the interactive committed-message issuance layer is one-more unforgeable if every PPT
adversary interacting in at most \(Q\) issuance sessions with the honest issuer outputs \(Q+1\) valid extracted
credential witnesses on pairwise distinct semantic messages only with negligible probability.  In this paper this
property is used as an explicit assumption in the ARC soundness decomposition.
\end{definition}

\begin{remark}[Scope of the issuance layer]\label{rem:bs-scope}
The protocol studied in Section~\ref{sec:blind-signature} is used as an issuer-view-hiding committed-message
issuance protocol together with a raw post-issuance verification relation.  One-more unforgeability remains a
separate assumption in the ARC soundness decomposition.
\end{remark}

\subsection{Anonymous Rate-Limiting Credentials}
\label{subsec:arc-security}
Rate limiting is part of the ARC syntax: the holder presents a credential together with a public nonce and a
PRF-derived public tag, while the verifier enforces the policy by maintaining state on previously accepted
\((\nonce,\prftag)\) pairs.

\begin{definition}[ARC syntax]\label{def:arc-syntax}
Let \(\nonceSpace\) be a finite nonce set.  An anonymous rate-limiting credential scheme is a tuple
\[
  \ARC=(\Setup,\IKeyGen,\Issue,\Show,\Verify)
\]
where \(\Issue\) outputs a credential, \(\Show(\credential,\nonce)\) outputs a presentation, and
\(\Verify(\vk,\presentation,\algstate)\) checks the presentation and updates verifier state.
\end{definition}

\begin{definition}[ARC soundness]\label{def:arc-soundness}
Let \(\ARC\) have nonce space \(\nonceSpace\).  In the ARC soundness experiment, the adversary receives public
parameters and an issuer verification key, interacts with the honest issuer in at most \(Q\) issuance sessions, and
then adaptively submits presentations to the verifier.  The verifier maintains its state according to the prescribed
policy and increments a counter \(N_{\mathsf{acc}}\) whenever a presentation is accepted.  Soundness requires
\[
  \Pr\bigl[N_{\mathsf{acc}}>Q|\nonceSpace|\bigr]\le\negl(\lambda)
\]
for every PPT adversary.
\end{definition}

\begin{definition}[Presentation unlinkability]\label{def:pres-unlink}
Presentation unlinkability is defined by a two-world experiment on distinct fresh nonces.  In world \(0\), the
challenger issues one credential and returns two presentations generated from that credential under the two nonces.
In world \(1\), it issues two independent credentials on the same admissible semantic message distribution and returns
one presentation from each credential.  The scheme has presentation unlinkability if every PPT distinguisher has
negligible advantage in distinguishing the two worlds, except for the intended linkage that follows from nonce reuse
or repeated public tags.
\end{definition}

\begin{remark}[Scope of the ARC theorems]\label{rem:arc-scope}
Section~\ref{sec:arc-construction} records conditional policy-soundness and presentation-unlinkability statements
for the construction induced by the committed-message issuance layer.  These statements rely on the explicitly
stated issuance assumption, NIZK knowledge soundness and zero knowledge, commitment binding, and PRF security.
\end{remark}

\subsection{Pseudorandom functions}
\label{subsec:prelim-prf}
\label{subsec:prf}

\begin{definition}[Pseudorandom function]\label{def:prf}
Let \(\mathcal K,\mathcal D,\mathcal T\) be finite sets.  A family
\[
  \PRF:\mathcal K\times \mathcal D\to \mathcal T
\]
is pseudorandom if oracle access to \(\PRF(k,\cdot)\) for uniform \(k\getsunif\mathcal K\) is computationally
indistinguishable from oracle access to a uniform random function from \(\mathcal D\) to \(\mathcal T\).
\end{definition}

We instantiate the PRF with a Poseidon2-style permutation, but the PRF security of the keyed
feed-forward/truncated family below is an explicit assumption rather than a theorem derived from the Poseidon2
permutation paper.

\begin{definition}[Poseidon2-style permutation]\label{def:poseidon-perm}
Let \(q\) be prime, let \(d\ge 3\) satisfy \(\gcd(d,q-1)=1\), and let \(R_F,R_P\in\mathbb N\) with \(R_F\) even.
A Poseidon2-style permutation alternates full external rounds and partial internal rounds, together with public
round constants and invertible linear mixing maps over \(\Fq\).
\end{definition}

\begin{definition}[Poseidon2-style PRF]\label{def:poseidon-prf}
For key \(\veck\in\mathcal K_{\mathsf{key}}\subseteq\Fq^{\ell_{\mathsf{key}}}\) and public input
\(\vecn\in\Fq^{\ell_{\mathsf{nonce}}}\), define
\[
  \PRF(\veck,\vecn)=\Tr_{\ell_{\mathsf{tag}}}\bigl(P(\veck\Vert\vecn)+(\veck\Vert\vecn)\bigr),
\]
where \(P\) is the Poseidon2-style permutation and \(\Tr_{\ell_{\mathsf{tag}}}\) returns the first
\(\ell_{\mathsf{tag}}\) lanes.
\end{definition}

\begin{assumption}[Concrete Poseidon2-style PRF assumption]\label{ass:poseidon-prf}
For the concrete parameter choices fixed in Appendix~\ref{app:prf}, the family
\[
  \PRF(k,n)=\Tr_{\ell_{\mathsf{tag}}}\bigl(P(k\Vert n)+(k\Vert n)\bigr)
\]
of Definition~\ref{def:poseidon-prf} is a secure pseudorandom function for
\(k\getsunif\mathcal K_{\mathsf{key}}\).  If the signed key space is changed, this assumption is made for the
changed key distribution and the ARC relations must enforce that same key space.
\end{assumption}
